\section*{Capítulo \ref{ch:g12:mechanical-waves} --- Ondas mecánicas}

\begin{solution}{exo:g12:mechanical-waves:1}
(a) $v = 320/40 = \qty{8.0}{m/s}$. (b) Ningún espectador recorre el
anillo: lo que viaja es la perturbación (el acto de ponerse de pie), no
el público. (c) $12/1.5 = \qty{8.0}{m/s}$ --- coherente.
\end{solution}

\begin{solution}{exo:g12:mechanical-waves:2}
(a) \omterm{def:g12:mechanical-waves:transverse}{transversal}; (b) \omterm{def:g12:mechanical-waves:transverse}{longitudinal}; (c) \omterm{def:g12:mechanical-waves:transverse}{longitudinal}; (d) \omterm{def:g12:mechanical-waves:transverse}{transversal};
(e) P \omterm{def:g12:mechanical-waves:transverse}{longitudinal}, S \omterm{def:g12:mechanical-waves:transverse}{transversal}.
\end{solution}

\begin{solution}{exo:g12:mechanical-waves:3}
(a) $d = 340 \times 4.0 \approx \qty{1.4}{km}$.
(b) $1500 \times 4.0 = \qty{6.0}{km}$. (c) La \omterm{def:g12:mechanical-waves:speed}{velocidad de propagación}
en el medio atravesado.
\end{solution}

\begin{solution}{exo:g12:mechanical-waves:4}
$\lambda = v/f$: $340/440 = \qty{0.77}{m}$ en el aire y
$1500/440 = \qty{3.4}{m}$ en el agua. La \omterm{def:g10:signals-and-waves:frequency}{frecuencia} se mantiene fija (la
dicta la fuente); $v$ y $\lambda$ cambian las dos.
\end{solution}

\begin{solution}{exo:g12:mechanical-waves:5}
$\lambda = \qty{12}{cm}$; $T = \lambda/v = 0.12/0.30 = \qty{0.40}{s}$;
$f = 1/T = \qty{2.5}{Hz}$.
\end{solution}

\begin{solution}{exo:g12:mechanical-waves:6}
(a) $v = 2.4/0.30 = \qty{8.0}{m/s}$. (b) Desde la primera foto le quedan
$\qty{6.0}{m}$ a la cresta: $t = 6.0/8.0 = \qty{0.75}{s}$ después de
ella. (c) La cinta sube y después baja, en dirección puramente
\omterm{def:g12:mechanical-waves:transverse}{transversal}, y se queda en $x = \qty{5.0}{m}$: la cuerda no viaja.
\end{solution}

\begin{solution}{exo:g12:mechanical-waves:7}
(a) Velocidad de acercamiento \qty{4.0}{m/s} y separación \qty{4.0}{m}:
$t = \qty{1.0}{s}$. (b) Superposición:
$3.0 + 2.0 = \qty{5.0}{cm}$. (c) $3.0 - 2.0 = \qty{1.0}{cm}$. (d) Salen
dos pulsos intactos que continúan sin cambios, como si no se hubieran
encontrado nunca.
\end{solution}

\begin{solution}{exo:g12:mechanical-waves:8}
(a) La fina: con la misma tensión, menos masa por \omterm{def:g10:orders-of-magnitude:unit}{metro} significa una
\omterm{def:g10:signals-and-waves:wave}{onda} más rápida. (b) Aumenta (más tensión). (c) No: la velocidad depende
únicamente del medio, no del tamaño del pulso.
\end{solution}

\begin{solution}{exo:g12:mechanical-waves:9}
Raíl: $1800/5900 \approx \qty{0.31}{s}$; aire:
$1800/340 \approx \qty{5.3}{s}$; intervalo $\approx \qty{5.0}{s}$.
\end{solution}

\begin{solution}{exo:g12:mechanical-waves:10}
Con $\lambda = \qty{24}{cm}$: (a) $2\lambda$, \omterm{def:g12:mechanical-waves:phase}{en fase}; (b)
$\lambda/2$, oposición; (c) $\tfrac32\lambda$, oposición; (d)
$1.25\lambda$, ninguna; (e) $\lambda/4$, ninguna.
\end{solution}

\begin{solution}{exo:g12:mechanical-waves:11}
$\lambda_P = 6.0 \times 2.0 = \qty{12}{km}$;
$\lambda_S = 3.5 \times 2.0 = \qty{7.0}{km}$ --- cientos de veces el
tamaño de un edificio. Los puntos mucho más próximos que $\lambda$ están
casi \omterm{def:g12:mechanical-waves:phase}{en fase}, así que el barrio entero se mueve a la vez.
\end{solution}

\begin{solution}{exo:g12:mechanical-waves:12}
(a) $T = \qty{0.50}{s}$, $f = \qty{2.0}{Hz}$.
(b) $\lambda = \qty{0.75}{m}$.
(c) $v = \lambda/T = 0.75/0.50 = \qty{1.5}{m/s}$.
(d) En $t = T = \qty{0.50}{s}$. (e) $1.875/0.75 = 2.5$ \omterm{def:g12:mechanical-waves:wavelength}{longitudes de
onda} $= (2 + \tfrac12)\lambda$: \omterm{def:g12:mechanical-waves:phase}{oposición de fase}.
\end{solution}

\begin{solution}{exo:g12:mechanical-waves:13}
(a) $T = 80/10 = \qty{8.0}{s}$, $f = \qty{0.125}{Hz}$.
(b) $v = \lambda/T = 120/8.0 = \qty{15}{m/s} = \qty{54}{km/h}$.
(c) $t = 900/15 = \qty{60}{s}$. (d) No: las \omterm{def:g10:signals-and-waves:wave}{ondas} no transportan materia
--- la boya fondeada solo cabecea.
\end{solution}

\begin{solution}{exo:g12:mechanical-waves:14}
(a) Se acercan a \qty{6.0}{m/s} a lo largo de \qty{6.0}{m}:
$t = \qty{1.0}{s}$; los desplazamientos opuestos se cancelan --- la
cuerda queda plana por un instante. (b) No: sus puntos se están moviendo
(las velocidades \omterm{def:g12:mechanical-waves:transverse}{transversales} se suman, no se cancelan); toda la
\omterm{def:g10:energy-conservation:energy}{energía} es cinética. (c) Los dos pulsos reaparecen y continúan sin
cambios. (d) La cuerda plana lleva velocidad y, por tanto, \omterm{def:g10:energy-conservation:energy}{energía}: la
\omterm{def:g10:signals-and-waves:wave}{onda} está guardada por un instante en el movimiento, no en la forma.
\end{solution}

\begin{solution}{exo:g12:mechanical-waves:15}
(a) $2\pi r$: \qty{3.1}{m} y \qty{50}{m} --- $\times 16$. (b) La \omterm{def:g10:energy-conservation:energy}{energía}
por \omterm{def:g10:orders-of-magnitude:unit}{metro} cae en un factor $16$, así que la \omterm{def:g10:signals-and-waves:amplitude}{amplitud} baja (la \omterm{def:g10:energy-conservation:energy}{energía}
crece con la \omterm{def:g10:signals-and-waves:amplitude}{amplitud}). (c) Los frentes rectos no se alargan: solo queda
la absorción. (d) El medio es activo: cada espectador se levanta con la
\omterm{def:g10:energy-conservation:energy}{energía} de sus propios músculos y realimenta la \omterm{def:g10:signals-and-waves:wave}{onda}.
\end{solution}

\begin{solution}{pb:g12:mechanical-waves:1}
\textbf{1.} Las P comprimen en el sentido de la marcha: \omterm{def:g12:mechanical-waves:transverse}{longitudinales};
las S sacuden de lado: \omterm{def:g12:mechanical-waves:transverse}{transversales}.

\textbf{2.} La P, más rápida, llega antes a todas partes.
$t_P = 120/6.0 = \qty{20}{s}$; $t_S = 120/3.5 \approx \qty{34}{s}$.

\textbf{3.} $\Delta t = d/v_S - d/v_P = d\,(1/v_S - 1/v_P)$; para
\qty{120}{km}: $34.3 - 20 \approx \qty{14}{s}$.

\textbf{4.} Despejando $d$:
$d = \dfrac{v_P v_S}{v_P - v_S}\,\Delta t
= \dfrac{6.0 \times 3.5}{2.5}\,\Delta t = \qty{8.4}{km}$ por segundo de
intervalo.

\textbf{5.} $\Delta t \propto d$: el registro de una sola estación da su
distancia al epicentro sin conocer la hora de origen.

\textbf{6.} $d_A = 8.4 \times 25 = \qty{2.1e2}{km}$.

\textbf{7.} El intervalo da una distancia, no una dirección: el
epicentro está en algún punto de la circunferencia de radio $d_A$
centrada en A.

\textbf{8.} $d_B = 8.4 \times 15 = \qty{1.3e2}{km}$;
$d_C = 8.4 \times 32 \approx \qty{2.7e2}{km}$.

\textbf{9.} Dos circunferencias se cortan en dos puntos; la tercera
selecciona uno: la intersección común es el epicentro.

\textbf{10.} $t_P = 210/6.0 = \qty{35}{s}$: el terremoto empezó a las
$09{:}14{:}32 - \qty{35}{s} = 09{:}13{:}57$.

\textbf{11.} $\qty{200}{m/s} = \qty{720}{km/h}$ --- un avión de línea en
crucero.

\textbf{12.} $t = \num{6.0e5}/200 = \qty{3.0e3}{s} = \qty{50}{min}$.

\textbf{13.} $t_P = 600/6.0 = \qty{100}{s}$. Aviso:
$3000 - 100 = \qty{2.9e3}{s} \approx \qty{48}{min}$.

\textbf{14.} $\lambda = vT = 200 \times 900 = \qty{1.8e5}{m} =
\qty{180}{km}$: un \omterm{def:g10:orders-of-magnitude:unit}{metro} escaso de subida repartido en \qty{180}{km}
--- el barco asciende una pendiente imperceptible a lo largo de varios
minutos.

\textbf{15.} $\lambda = 10 \times 900 = \qty{9.0}{km}$: la \omterm{def:g10:signals-and-waves:wave}{onda} se
apretuja veinte veces, su \omterm{def:g10:energy-conservation:energy}{energía} se amontona en una longitud más corta
y, por tanto, su altura tiene que crecer --- el tsunami se encabrita al
llegar a la orilla.

\textbf{16.} $2\pi r$: \qty{63}{km} y después
$\approx \qty{1.3e3}{km}$ --- $\times 21$.

\textbf{17.} La misma \omterm{def:g10:energy-conservation:energy}{energía} se reparte sobre un frente $21$ veces más
largo, así que la \omterm{def:g10:signals-and-waves:amplitude}{amplitud} baja incluso en una corteza sin pérdidas; una
corteza real añade la absorción.

\textbf{18.} Las \omterm{def:g10:signals-and-waves:wave}{ondas} S, cizalla \omterm{def:g12:mechanical-waves:transverse}{transversal}, no pueden atravesar un
líquido: la sombra de las \omterm{def:g10:signals-and-waves:wave}{ondas} S demuestra que el \omterm{def:g11:fundamental-interactions:ladder}{núcleo} (externo) es
líquido.

\textbf{19.} $\num{1e15}/\num{1.5e9} \approx \num{6.7e5}$: unos
setecientos mil trenes a toda velocidad.

\textbf{20.} El terremoto ocurrió a las $09{:}13{:}57$, mar adentro en
el cruce de las tres circunferencias, a \qty{210}{km} de la estación A.
La alarma de la \omterm{def:g10:signals-and-waves:wave}{onda} P del puerto sonó \qty{100}{s} después del origen,
unos \qty{48}{min} antes de que llegase el tsunami.
\end{solution}
